\documentclass[9pt,a4paper]{article}
\usepackage[margin=0.55in]{geometry}
\usepackage{booktabs}
\usepackage{enumitem}
\usepackage[hidelinks]{hyperref}
\usepackage{microtype}
\setlength{\parindent}{0pt}
\setlength{\parskip}{2.5pt}
\setlist{nosep,leftmargin=1.2em}
\newcommand{\sect}[1]{\vspace{2pt}\textbf{#1}\quad}
\pagestyle{empty}

\begin{document}
\begin{center}
{\large\bfseries Causal Dialog Route Reranker}\\[-1pt]
{\small SE team \quad Alternative DialNav Challenge entry, August 2026}
\end{center}
\vspace{-5pt}

\sect{Method and communication boundary.}
We submit a single-rollout dialog-navigation method built on the official
RAINbow system. Only the Navigator owns and invokes the frozen LANA question
generator (QG). At a dialog step, the Navigator describes its current
observation as a natural-language question. The Guide, which has no QG module,
localizes that question with frozen GTL and generates a natural-language route
answer with the frozen LANA answer generator. Only the question and answer
strings enter the Navigator's instruction state.

Our addition is Guide-private temporal localization. If the previous Guide
answer was generated from route $R$ at step $t_0$, the expected location at
step $t$ is $e_t=R_{\min(\max(t-t_0,0),|R|-1)}$. For GTL top-five candidate
$v$ with probability $p(v)$, the Guide selects
$\arg\min_v[-\log(\max(p(v),10^{-12}))+0.5d_G(v,e_t)]$. Without prior route
state, it uses GTL top-1. The state never crosses to the Navigator as structured
data. A Guide goal confirmation is also natural-language text. The system
executes one causal trajectory, with no replay or episode selector.

This alternative uses no shared QG, question fingerprint, target-language
signature, target-description transfer, language DFS, visual signature,
trajectory splicing, evaluation cache/label, manual evaluation annotation,
Gemini, UniAPI, or hosted inference API.

\sect{Data, models, selection, and execution.}
The base is RAINbow commit \texttt{4ef165fad77e67502695}\\[-1pt]
\texttt{8e16d9e516e523192a33}. Frozen components are the
official DST Navigator, LANA QG, LANA answer generator, and GTL localizer; no
weights were trained or fine-tuned. Their upstream data provenance is
RAIN/RxR/R2R/CVDN and Matterport3D. No additional dataset was used.

Top-$k=5$ and weight $0.5$ were fixed using 806 train episodes from 30
scan-disjoint scans. Relative to the same-fold base, train Score increased from
60.2713 to 64.2928 (+4.0215 points), with scan-clustered 95\% CI
[+2.1385,+6.0464]. Final inference used seed 0, one NVIDIA H100 GPU, Python
3.10, PyTorch 2.10, batch size 8, WTA \texttt{ctc\_0.7\_2}, answer-path limit
20, and a 50-action cap.

\sect{Results and limitations.}
\begin{center}
\small
\begin{tabular}{lrrrrr}
\toprule
Split & Episodes & Score & SR & Steps & Dialogs \\
\midrule
val-seen & 91 & 0.5977 & 65.93\% & 21.59 & 3.26 \\
val-unseen & 241 & 0.2646 & 31.95\% & 24.59 & 5.29 \\
released Test & 285 & hidden & 27.37\% & 27.37 & 6.05 \\
\bottomrule
\end{tabular}
\end{center}
\vspace{-3pt}
The released-Test result is 78/285 and is not the hidden ranking Test Score.
The method was frozen from train-only evidence before its saved Test run; the
official outcomes did not alter its parameters. Its main limitation is weak
unseen localization: an incorrect Guide localization can still produce an
incorrect route or confirmation.

\sect{Reproducibility.}
The delivery includes the source overlay, fixed configuration, tests, exact
submission, checksums, model/data card, and runtime-weight hashes. Submission
SHA256 is \texttt{52c191922ff14423efe5b290c0a9e1aa}\\[-1pt]
\texttt{0fa0e094c06227145db3a421fbbf5efc}.

\end{document}
